\documentclass[a4paper,10pt]{article}
\usepackage{eumat}

\begin{document}
\begin{eulernotebook}
\eulerheading{Menggambar Plot 3D dengan EMT}
\begin{eulercomment}
Ini adalah pengantar plot 3D di Euler. Kita memerlukan plot 3D untuk memvisualisasikan fungsi dua variabel.

Euler menggambar fungsi semacam itu menggunakan algoritma pengurutan untuk menyembunyikan bagian di latar belakang.
Pada umumnya, Euler menggunakan proyeksi sentral. Defaultnya adalah dari kuadran x-y positif menuju titik asal
x=y=z=0, tetapi angle=0° memandang searah sumbu y. Sudut pandang dan ketinggian dapat diubah.

Euler dapat menggambar

- permukaan dengan arsiran (shading) dan garis level atau rentang level,\\
- awan titik,\\
- kurva parametrik,\\
- permukaan implisit.

Plot 3D suatu fungsi menggunakan plot3d. Cara termudah adalah menggambar ekspresi dalam x dan y. Parameter r
mengatur rentang plot di sekitar (0,0).
\end{eulercomment}
\begin{eulerprompt}
>aspect(1.5); plot3d("x^2+sin(y)",-5,5,0,6*pi):
\end{eulerprompt}
\eulerimg{17}{images/EMT4Plot3D_Indonesia-001.png}
\begin{eulerprompt}
>plot3d("x^2+x*sin(y)",-5,5,0,6*pi):
\end{eulerprompt}
\eulerimg{17}{images/EMT4Plot3D_Indonesia-002.png}
\begin{eulercomment}
Silakan lakukan modifikasi agar gambar "talang bergelombang" tersebut tidak lurus melainkan melengkung/melingkar, baik
melingkar secara mendatar maupun melingkar turun/naik (seperti papan peluncur pada kolam renang. Temukan rumusnya.
\end{eulercomment}
\eulerheading{Fungsi dengan Dua Variabel}
\begin{eulercomment}
Untuk grafik suatu fungsi, gunakan

- ekspresi sederhana dalam x dan y,\\
- nama fungsi dua variabel,\\
- atau matriks data.

Defaultnya adalah kisi kawat terisi dengan warna berbeda pada kedua sisi. Perhatikan bahwa banyaknya interval grid
default adalah 10, tetapi plot menggunakan 40x40 persegi panjang default untuk membangun permukaan. Hal ini dapat
diubah.

- n=40, n=[40,40]: banyaknya garis grid pada setiap arah\\
- grid=10, grid=[10,10]: banyaknya garis grid pada setiap arah.

Kita gunakan default n=40 dan grid=10.
\end{eulercomment}
\begin{eulerprompt}
>plot3d("x^2+y^2"):
\end{eulerprompt}
\eulerimg{17}{images/EMT4Plot3D_Indonesia-003.png}
\begin{eulercomment}
Interaksi pengguna dimungkinkan dengan parameter \textgreater{}user. Pengguna dapat menekan tombol-tombol berikut.

- kiri,kanan,atas,bawah: memutar sudut pandang\\
- +,-: memperbesar atau memperkecil\\
- a: membuat anaglyph (lihat di bawah)\\
- l: mengaktifkan/menonaktifkan pemutaran sumber cahaya (lihat di bawah)\\
- spasi: kembali ke default\\
- return: mengakhiri interaksi
\end{eulercomment}
\begin{eulerprompt}
>plot3d("exp(-x^2+y^2)",>user, ...
>  title="Turn with the vector keys (press return to finish)"):
\end{eulerprompt}
\eulerimg{17}{images/EMT4Plot3D_Indonesia-004.png}
\begin{eulercomment}
Rentang plot untuk fungsi dapat ditentukan dengan

- a,b: rentang x\\
- c,d: rentang y\\
- r: persegi simetris di sekitar (0,0).\\
- n: banyaknya subinterval untuk plot.

Ada beberapa parameter untuk menskalakan fungsi atau mengubah tampilan grafik.

fscale: menskalakan ke nilai fungsi (defaultnya \textless{}fscale).\\
scale: bilangan atau vektor 1x2 untuk menskalakan ke arah x dan y.\\
frame: jenis bingkai (default 1).
\end{eulercomment}
\begin{eulerprompt}
>plot3d("exp(-(x^2+y^2)/5)",r=10,n=80,fscale=4,scale=1.2,frame=3,>user):
\end{eulerprompt}
\eulerimg{17}{images/EMT4Plot3D_Indonesia-005.png}
\begin{eulercomment}
Tampilan dapat diubah dengan banyak cara.

- distance: jarak pandang ke plot.\\
- zoom: nilai zoom.\\
- angle: sudut terhadap sumbu y negatif dalam radian.\\
- height: ketinggian pandangan dalam radian.

Nilai default dapat diperiksa atau diubah dengan fungsi view(). Fungsi ini mengembalikan parameter dengan urutan di
atas.
\end{eulercomment}
\begin{eulerprompt}
>view
\end{eulerprompt}
\begin{euleroutput}
  [5,  2.6,  2,  0.4]
\end{euleroutput}
\begin{eulercomment}
Jarak yang lebih dekat memerlukan zoom yang lebih kecil. Efeknya lebih mirip lensa sudut lebar.

Pada contoh berikut, angle=0 dan height=0 memandang dari sumbu y negatif. Label sumbu y disembunyikan dalam kasus
ini.
\end{eulercomment}
\begin{eulerprompt}
>plot3d("x^2+y",distance=3,zoom=1,angle=pi/2,height=0):
\end{eulerprompt}
\eulerimg{17}{images/EMT4Plot3D_Indonesia-006.png}
\begin{eulercomment}
Plot selalu memandang ke pusat kubus plot. Anda dapat memindahkan pusat tersebut dengan parameter center.
\end{eulercomment}
\begin{eulerprompt}
>plot3d("x^4+y^2",a=0,b=1,c=-1,d=1,angle=-20°,height=20°, ...
>  center=[0.4,0,0],zoom=5):
\end{eulerprompt}
\eulerimg{17}{images/EMT4Plot3D_Indonesia-007.png}
\begin{eulercomment}
Plot diskalakan agar muat dalam kubus satuan untuk dilihat. Jadi tidak perlu mengubah jarak atau zoom bergantung
pada ukuran plot. Namun, label merujuk pada ukuran sebenarnya.

Jika Anda mematikannya dengan scale=false, Anda perlu memastikan plot tetap muat dalam jendela plot, dengan mengubah
jarak pandang atau zoom, dan memindahkan pusat.
\end{eulercomment}
\begin{eulerprompt}
>plot3d("5*exp(-x^2-y^2)",r=2,<fscale,<scale,distance=13,height=50°, ...
>  center=[0,0,-2],frame=3):
\end{eulerprompt}
\eulerimg{17}{images/EMT4Plot3D_Indonesia-008.png}
\begin{eulercomment}
Plot polar juga tersedia. Parameter polar=true menggambar plot polar. Fungsinya tetap harus berupa fungsi x dan y.
Parameter "fscale" menskalakan fungsi dengan skala tersendiri. Jika tidak, fungsi diskalakan agar muat dalam kubus.
\end{eulercomment}
\begin{eulerprompt}
>plot3d("1/(x^2+y^2+1)",r=5,>polar, ...
>fscale=2,>hue,n=100,zoom=4,>contour,color=blue):
\end{eulerprompt}
\eulerimg{17}{images/EMT4Plot3D_Indonesia-009.png}
\begin{eulerprompt}
>function f(r) := exp(-r/2)*cos(r); ...
>plot3d("f(x^2+y^2)",>polar,scale=[1,1,0.4],r=pi,frame=3,zoom=4):
\end{eulerprompt}
\eulerimg{17}{images/EMT4Plot3D_Indonesia-010.png}
\begin{eulercomment}
Parameter rotate memutar fungsi dalam x mengelilingi sumbu x.

- rotate=1: Menggunakan sumbu x\\
- rotate=2: Menggunakan sumbu z
\end{eulercomment}
\begin{eulerprompt}
>plot3d("x^2+1",a=-1,b=1,rotate=true,grid=5):
\end{eulerprompt}
\eulerimg{17}{images/EMT4Plot3D_Indonesia-011.png}
\begin{eulerprompt}
>plot3d("x^2+1",a=-1,b=1,rotate=2,grid=5):
\end{eulerprompt}
\eulerimg{17}{images/EMT4Plot3D_Indonesia-012.png}
\begin{eulerprompt}
>plot3d("sqrt(25-x^2)",a=0,b=5,rotate=1):
\end{eulerprompt}
\eulerimg{17}{images/EMT4Plot3D_Indonesia-013.png}
\begin{eulerprompt}
>plot3d("x*sin(x)",a=0,b=6pi,rotate=2):
\end{eulerprompt}
\eulerimg{17}{images/EMT4Plot3D_Indonesia-014.png}
\begin{eulercomment}
Berikut plot dengan tiga fungsi.
\end{eulercomment}
\begin{eulerprompt}
>plot3d("x","x^2+y^2","y",r=2,zoom=3.5,frame=3):
\end{eulerprompt}
\eulerimg{17}{images/EMT4Plot3D_Indonesia-015.png}
\eulerheading{Plot Kontur}
\begin{eulercomment}
Untuk plot, Euler menambahkan garis grid. Sebagai gantinya, dapat digunakan garis level dan rona (hue) satu warna
atau rona berwarna spektral. Euler dapat menggambar ketinggian fungsi pada plot dengan arsiran. Pada semua plot 3D,
Euler dapat menghasilkan anaglyph merah/cyan.

- \textgreater{}hue: Menyalakan arsiran ringan sebagai ganti kawat.\\
- \textgreater{}contour: Menggambar garis kontur otomatis pada plot.\\
- level=... (atau levels): Vektor nilai untuk garis kontur.

Defaultnya adalah level="auto", yang menghitung beberapa garis level secara otomatis. Seperti terlihat pada plot,
level tersebut sebenarnya adalah rentang level.

Gaya default dapat diubah. Untuk plot kontur berikut, kita menggunakan grid yang lebih halus yaitu 100x100 titik,
menskalakan fungsi dan plot, dan menggunakan sudut pandang yang berbeda.
\end{eulercomment}
\begin{eulerprompt}
>plot3d("exp(-x^2-y^2)",r=2,n=100,level="thin", ...
> >contour,>spectral,fscale=1,scale=1.1,angle=45°,height=20°):
\end{eulerprompt}
\eulerimg{17}{images/EMT4Plot3D_Indonesia-016.png}
\begin{eulerprompt}
>plot3d("exp(x*y)",angle=100°,>contour,color=green):
\end{eulerprompt}
\eulerimg{17}{images/EMT4Plot3D_Indonesia-017.png}
\begin{eulercomment}
Arsiran default menggunakan warna abu-abu. Tetapi rentang warna spektral juga tersedia.

- \textgreater{}spectral: Menggunakan skema spektral default\\
- color=...: Menggunakan warna khusus atau skema spektral

Untuk plot berikut, kita menggunakan skema spektral default dan menambah banyaknya titik agar tampilannya sangat
mulus.
\end{eulercomment}
\begin{eulerprompt}
>plot3d("x^2+y^2",>spectral,>contour,n=100):
\end{eulerprompt}
\eulerimg{17}{images/EMT4Plot3D_Indonesia-018.png}
\begin{eulercomment}
Sebagai ganti garis level otomatis, kita juga dapat menetapkan nilai garis level. Ini akan menghasilkan garis level
tipis, bukan rentang level.
\end{eulercomment}
\begin{eulerprompt}
>plot3d("x^2-y^2",0,5,0,5,level=-1:0.1:1,color=redgreen):
\end{eulerprompt}
\eulerimg{17}{images/EMT4Plot3D_Indonesia-019.png}
\begin{eulercomment}
Pada plot berikut, kita menggunakan dua pita level yang sangat lebar dari -0.1 sampai 1, dan dari 0.9 sampai 1. Ini
dimasukkan sebagai matriks dengan batas level sebagai kolom.

Selain itu, kita menumpangkan grid dengan 10 interval pada setiap arah.
\end{eulercomment}
\begin{eulerprompt}
>plot3d("x^2+y^3",level=[-0.1,0.9;0,1], ...
>  >spectral,angle=30°,grid=10,contourcolor=gray):
\end{eulerprompt}
\eulerimg{17}{images/EMT4Plot3D_Indonesia-020.png}
\begin{eulercomment}
Pada contoh berikut, kita menggambar himpunan di mana

\end{eulercomment}
\begin{eulerformula}
\[
f(x,y) = x^y-y^x = 0
\]
\end{eulerformula}
\begin{eulercomment}
Kita menggunakan satu garis tipis untuk garis level.
\end{eulercomment}
\begin{eulerprompt}
>plot3d("x^y-y^x",level=0,a=0,b=6,c=0,d=6,contourcolor=red,n=100):
\end{eulerprompt}
\eulerimg{17}{images/EMT4Plot3D_Indonesia-021.png}
\begin{eulercomment}
Dimungkinkan menampilkan bidang kontur di bawah plot. Warna dan jarak ke plot dapat ditentukan.
\end{eulercomment}
\begin{eulerprompt}
>plot3d("x^2+y^4",>cp,cpcolor=green,cpdelta=0.2):
\end{eulerprompt}
\eulerimg{17}{images/EMT4Plot3D_Indonesia-022.png}
\begin{eulercomment}
Berikut beberapa gaya lainnya. Kita selalu mematikan bingkai, dan menggunakan berbagai skema warna untuk plot dan
grid.
\end{eulercomment}
\begin{eulerprompt}
>figure(2,2); ...
>expr="y^3-x^2"; ...
>figure(1);  ...
>  plot3d(expr,<frame,>cp,cpcolor=spectral); ...
>figure(2);  ...
>  plot3d(expr,<frame,>spectral,grid=10,cp=2); ...
>figure(3);  ...
>  plot3d(expr,<frame,>contour,color=gray,nc=5,cp=3,cpcolor=greenred); ...
>figure(4);  ...
>  plot3d(expr,<frame,>hue,grid=10,>transparent,>cp,cpcolor=gray); ...
>figure(0):
\end{eulerprompt}
\eulerimg{17}{images/EMT4Plot3D_Indonesia-023.png}
\begin{eulercomment}
Ada beberapa skema spektral lain, bernomor 1 sampai 9. Tetapi Anda juga dapat menggunakan color=nilai, dengan nilai
berupa

- spectral: untuk rentang dari biru ke merah\\
- white: untuk rentang yang lebih pudar\\
- yellowblue,purplegreen,blueyellow,greenred\\
- blueyellow, greenpurple,yellowblue,redgreen
\end{eulercomment}
\begin{eulerprompt}
>figure(3,3); ...
>for i=1:9;  ...
>  figure(i); plot3d("x^2+y^2",spectral=i,>contour,>cp,<frame,zoom=4);  ...
>end; ...
>figure(0):
\end{eulerprompt}
\eulerimg{17}{images/EMT4Plot3D_Indonesia-024.png}
\begin{eulercomment}
Sumber cahaya dapat diubah dengan l dan tombol kursor selama interaksi pengguna. Sumber cahaya juga dapat diatur
dengan parameter.

- light: arah cahaya\\
- amb: cahaya sekitar (ambient) antara 0 dan 1

Perhatikan bahwa program tidak membedakan sisi-sisi plot. Tidak ada bayangan. Untuk itu Anda memerlukan Povray.
\end{eulercomment}
\begin{eulerprompt}
>plot3d("-x^2-y^2", ...
>  hue=true,light=[0,1,1],amb=0,user=true, ...
>  title="Press l and cursor keys (return to exit)"):
\end{eulerprompt}
\eulerimg{17}{images/EMT4Plot3D_Indonesia-025.png}
\begin{eulercomment}
Parameter color mengubah warna permukaan. Warna garis level juga dapat diubah.
\end{eulercomment}
\begin{eulerprompt}
>plot3d("-x^2-y^2",color=rgb(0.2,0.2,0),hue=true,frame=false, ...
>  zoom=3,contourcolor=red,level=-2:0.1:1,dl=0.01):
\end{eulerprompt}
\eulerimg{17}{images/EMT4Plot3D_Indonesia-026.png}
\begin{eulercomment}
Warna 0 memberikan efek pelangi khusus.
\end{eulercomment}
\begin{eulerprompt}
>plot3d("x^2/(x^2+y^2+1)",color=0,hue=true,grid=10):
\end{eulerprompt}
\eulerimg{17}{images/EMT4Plot3D_Indonesia-027.png}
\begin{eulercomment}
Permukaan juga dapat dibuat transparan.
\end{eulercomment}
\begin{eulerprompt}
>plot3d("x^2+y^2",>transparent,grid=10,wirecolor=red):
\end{eulerprompt}
\eulerimg{17}{images/EMT4Plot3D_Indonesia-028.png}
\eulerheading{Plot Implisit}
\begin{eulercomment}
Ada juga plot implisit dalam tiga dimensi. Euler membuat potongan melalui objek. Fitur plot3d mencakup plot
implisit. Plot ini menunjukkan himpunan nol suatu fungsi tiga variabel. Solusi dari

\end{eulercomment}
\begin{eulerformula}
\[
f(x,y,z) = 0
\]
\end{eulerformula}
\begin{eulercomment}
dapat divisualisasikan dalam potongan sejajar bidang x-y, x-z, dan y-z.

- implicit=1: potongan sejajar bidang y-z\\
- implicit=2: potongan sejajar bidang x-z\\
- implicit=4: potongan sejajar bidang x-y

Jumlahkan nilai-nilai ini jika Anda mau. Pada contoh kita menggambar

\end{eulercomment}
\begin{eulerformula}
\[
M = \{ (x,y,z) : x^2+y^3+zy=1 \}
\]
\end{eulerformula}
\begin{eulerprompt}
>plot3d("x^2+y^3+z*y-1",r=5,implicit=3):
\end{eulerprompt}
\eulerimg{17}{images/EMT4Plot3D_Indonesia-029.png}
\begin{eulerprompt}
>c=1; d=1;
>plot3d("((x^2+y^2-c^2)^2+(z^2-1)^2)*((y^2+z^2-c^2)^2+(x^2-1)^2)*((z^2+x^2-c^2)^2+(y^2-1)^2)-d",r=2,<frame,>implicit,>user): 
\end{eulerprompt}
\eulerimg{17}{images/EMT4Plot3D_Indonesia-030.png}
\begin{eulerprompt}
>plot3d("x^2+y^2+4*x*z+z^3",>implicit,r=2,zoom=2.5):
\end{eulerprompt}
\eulerimg{17}{images/EMT4Plot3D_Indonesia-031.png}
\eulerheading{Menggambar Data 3D}
\begin{eulercomment}
Sama seperti plot2d, plot3d menerima data. Untuk objek 3D, Anda perlu menyediakan matriks nilai x, y, dan z, atau
tiga fungsi atau ekspresi fx(x,y), fy(x,y), fz(x,y).

\end{eulercomment}
\begin{eulerformula}
\[
\gamma(t,s) = (x(t,s),y(t,s),z(t,s))
\]
\end{eulerformula}
\begin{eulercomment}
Karena x,y,z adalah matriks, kita mengasumsikan (t,s) menjelajahi grid persegi. Akibatnya, Anda dapat menggambar
bayangan persegi panjang di ruang.

Anda dapat menggunakan bahasa matriks Euler untuk menghasilkan koordinat secara efektif.

Pada contoh berikut, kita menggunakan vektor nilai t dan vektor kolom nilai s untuk memparameterkan permukaan bola.
Pada gambar kita dapat menandai daerah, dalam kasus kita daerah kutub.
\end{eulercomment}
\begin{eulerprompt}
>t=linspace(0,2pi,180); s=linspace(-pi/2,pi/2,90)'; ...
>x=cos(s)*cos(t); y=cos(s)*sin(t); z=sin(s); ...
>plot3d(x,y,z,>hue, ...
>color=blue,<frame,grid=[10,20], ...
>values=s,contourcolor=red,level=[90°-24°;90°-22°], ...
>scale=1.4,height=50°):
\end{eulerprompt}
\eulerimg{17}{images/EMT4Plot3D_Indonesia-032.png}
\begin{eulercomment}
Berikut contoh yang merupakan grafik suatu fungsi.
\end{eulercomment}
\begin{eulerprompt}
>t=-1:0.1:1; s=(-1:0.1:1)'; plot3d(t,s,t*s,grid=10):
\end{eulerprompt}
\eulerimg{17}{images/EMT4Plot3D_Indonesia-033.png}
\begin{eulercomment}
Namun, kita dapat membuat berbagai macam permukaan. Berikut permukaan yang sama sebagai fungsi

\end{eulercomment}
\begin{eulerformula}
\[
x = y \, z
\]
\end{eulerformula}
\begin{eulerprompt}
>plot3d(t*s,t,s,angle=180°,grid=10):
\end{eulerprompt}
\eulerimg{17}{images/EMT4Plot3D_Indonesia-034.png}
\begin{eulercomment}
Dengan usaha lebih, kita dapat membuat banyak permukaan.

Pada contoh berikut kita membuat tampilan berarsir dari bola yang terdistorsi. Koordinat yang biasa untuk bola
adalah

\end{eulercomment}
\begin{eulerformula}
\[
\gamma(t,s) = (\cos(t)\cos(s),\sin(t)\sin(s),\cos(s))
\]
\end{eulerformula}
\begin{eulercomment}
dengan

\end{eulercomment}
\begin{eulerformula}
\[
0 \le t \le 2\pi, \quad \frac{-\pi}{2} \le s \le \frac{\pi}{2}.
\]
\end{eulerformula}
\begin{eulercomment}
Kita mendistorsi ini dengan faktor

\end{eulercomment}
\begin{eulerformula}
\[
d(t,s) = \frac{\cos(4t)+\cos(8s)}{4}.
\]
\end{eulerformula}
\begin{eulerprompt}
>t=linspace(0,2pi,320); s=linspace(-pi/2,pi/2,160)'; ...
>d=1+0.2*(cos(4*t)+cos(8*s)); ...
>plot3d(cos(t)*cos(s)*d,sin(t)*cos(s)*d,sin(s)*d,hue=1, ...
>  light=[1,0,1],frame=0,zoom=5):
\end{eulerprompt}
\eulerimg{17}{images/EMT4Plot3D_Indonesia-035.png}
\begin{eulercomment}
Tentu saja, awan titik juga dimungkinkan. Untuk menggambar data titik di ruang, kita memerlukan tiga vektor untuk
koordinat titik-titik tersebut.

Gaya-gayanya sama seperti pada plot2d dengan points=true;
\end{eulercomment}
\begin{eulerprompt}
>n=500;  ...
>  plot3d(normal(1,n),normal(1,n),normal(1,n),points=true,style="."):
\end{eulerprompt}
\eulerimg{17}{images/EMT4Plot3D_Indonesia-036.png}
\begin{eulercomment}
Dimungkinkan juga menggambar kurva dalam 3D. Dalam hal ini, lebih mudah menghitung titik-titik kurva terlebih
dahulu. Untuk kurva pada bidang kita menggunakan barisan koordinat dan parameter wire=true.
\end{eulercomment}
\begin{eulerprompt}
>t=linspace(0,8pi,500); ...
>plot3d(sin(t),cos(t),t/10,>wire,zoom=3):
\end{eulerprompt}
\eulerimg{17}{images/EMT4Plot3D_Indonesia-037.png}
\begin{eulerprompt}
>t=linspace(0,4pi,1000); plot3d(cos(t),sin(t),t/2pi,>wire, ...
>linewidth=3,wirecolor=blue):
\end{eulerprompt}
\eulerimg{17}{images/EMT4Plot3D_Indonesia-038.png}
\begin{eulerprompt}
>X=cumsum(normal(3,100)); ...
> plot3d(X[1],X[2],X[3],>anaglyph,>wire):
\end{eulerprompt}
\eulerimg{17}{images/EMT4Plot3D_Indonesia-039.png}
\begin{eulercomment}
EMT juga dapat menggambar dalam mode anaglyph. Untuk melihat plot seperti itu, Anda memerlukan kacamata merah/cyan.
\end{eulercomment}
\begin{eulerprompt}
> plot3d("x^2+y^3",>anaglyph,>contour,angle=30°):
\end{eulerprompt}
\eulerimg{17}{images/EMT4Plot3D_Indonesia-040.png}
\begin{eulercomment}
Sering kali, skema warna spektral digunakan untuk plot. Ini menonjolkan ketinggian fungsi.
\end{eulercomment}
\begin{eulerprompt}
>plot3d("x^2*y^3-y",>spectral,>contour,zoom=3.2):
\end{eulerprompt}
\eulerimg{17}{images/EMT4Plot3D_Indonesia-041.png}
\begin{eulercomment}
Euler juga dapat menggambar permukaan terparameter, ketika parameternya adalah nilai x, y, dan z dari bayangan grid
persegi panjang di ruang.

Untuk demo berikut, kita menyiapkan parameter u dan v, dan membangkitkan koordinat ruang darinya.
\end{eulercomment}
\begin{eulerprompt}
>u=linspace(-1,1,10); v=linspace(0,2*pi,50)'; ...
>X=(3+u*cos(v/2))*cos(v); Y=(3+u*cos(v/2))*sin(v); Z=u*sin(v/2); ...
>plot3d(X,Y,Z,>anaglyph,<frame,>wire,scale=2.3):
\end{eulerprompt}
\eulerimg{17}{images/EMT4Plot3D_Indonesia-042.png}
\begin{eulercomment}
Berikut contoh yang lebih rumit, yang tampak megah dengan kacamata merah/cyan.
\end{eulercomment}
\begin{eulerprompt}
>u:=linspace(-pi,pi,160); v:=linspace(-pi,pi,400)';  ...
>x:=(4*(1+.25*sin(3*v))+cos(u))*cos(2*v); ...
>y:=(4*(1+.25*sin(3*v))+cos(u))*sin(2*v); ...
> z=sin(u)+2*cos(3*v); ...
>plot3d(x,y,z,frame=0,scale=1.5,hue=1,light=[1,0,-1],zoom=2.8,>anaglyph):
\end{eulerprompt}
\eulerimg{17}{images/EMT4Plot3D_Indonesia-043.png}
\eulerheading{Plot Statistik}
\begin{eulercomment}
Plot batang juga dimungkinkan. Untuk itu, kita harus menyediakan

- x: vektor baris dengan n+1 elemen\\
- y: vektor kolom dengan n+1 elemen\\
- z: matriks nilai nxn.

z boleh lebih besar, tetapi hanya nilai nxn yang akan digunakan.

Pada contoh, kita pertama-tama menghitung nilai-nilainya. Kemudian kita menyesuaikan x dan y, sehingga vektor-vektor
tersebut berpusat pada nilai yang digunakan.
\end{eulercomment}
\begin{eulerprompt}
>x=-1:0.1:1; y=x'; z=x^2+y^2; ...
>xa=(x|1.1)-0.05; ya=(y_1.1)-0.05; ...
>plot3d(xa,ya,z,bar=true):
\end{eulerprompt}
\eulerimg{17}{images/EMT4Plot3D_Indonesia-044.png}
\begin{eulercomment}
Dimungkinkan membagi plot suatu permukaan menjadi dua bagian atau lebih.
\end{eulercomment}
\begin{eulerprompt}
>x=-1:0.1:1; y=x'; z=x+y; d=zeros(size(x)); ...
>plot3d(x,y,z,disconnect=2:2:20):
\end{eulerprompt}
\eulerimg{17}{images/EMT4Plot3D_Indonesia-045.png}
\begin{eulercomment}
Jika Anda memuat atau membangkitkan matriks data M dari sebuah file dan perlu menggambarnya dalam 3D, Anda dapat
menskalakan matriks ke [-1,1] dengan scale(M), atau menskalakan matriks dengan \textgreater{}zscale. Ini dapat digabungkan dengan
faktor skala individual yang diterapkan tambahan.
\end{eulercomment}
\begin{eulerprompt}
>i=1:20; j=i'; ...
>plot3d(i*j^2+100*normal(20,20),>zscale,scale=[1,1,1.5],angle=-40°,zoom=1.8):
\end{eulerprompt}
\eulerimg{17}{images/EMT4Plot3D_Indonesia-046.png}
\begin{eulerprompt}
>Z=intrandom(5,100,6); v=zeros(5,6); ...
>loop 1 to 5; v[#]=getmultiplicities(1:6,Z[#]); end; ...
>columnsplot3d(v',scols=1:5,ccols=[1:5]):
\end{eulerprompt}
\eulerimg{17}{images/EMT4Plot3D_Indonesia-047.png}
\eulerheading{Permukaan Benda Putar}
\begin{eulerprompt}
>plot2d("(x^2+y^2-1)^3-x^2*y^3",r=1.3, ...
>style="#",color=red,<outline, ...
>level=[-2;0],n=100):
\end{eulerprompt}
\eulerimg{17}{images/EMT4Plot3D_Indonesia-048.png}
\begin{eulerprompt}
>ekspresi &= (x^2+y^2-1)^3-x^2*y^3; $ekspresi
\end{eulerprompt}
\begin{eulerformula}
\[
\left(y^2+x^2-1\right)^3-x^2\,y^3
\]
\end{eulerformula}
\begin{eulercomment}
Kita ingin memutar kurva hati mengelilingi sumbu y. Berikut ekspresi yang mendefinisikan hati tersebut:

\end{eulercomment}
\begin{eulerformula}
\[
f(x,y)=(x^2+y^2-1)^3-x^2.y^3.
\]
\end{eulerformula}
\begin{eulercomment}
Selanjutnya kita tetapkan

\end{eulercomment}
\begin{eulerformula}
\[
x=r.cos(a),\quad y=r.sin(a).
\]
\end{eulerformula}
\begin{eulerprompt}
>function fr(r,a) &= ekspresi with [x=r*cos(a),y=r*sin(a)] | trigreduce; $fr(r,a)
\end{eulerprompt}
\begin{eulerformula}
\[
\left(r^2-1\right)^3+\frac{\left(\sin \left(5\,a\right)-\sin \left(
 3\,a\right)-2\,\sin a\right)\,r^5}{16}
\]
\end{eulerformula}
\begin{eulercomment}
Ini memungkinkan kita mendefinisikan fungsi numerik, yang menyelesaikan r jika a diberikan. Dengan fungsi tersebut
kita dapat menggambar hati yang diputar sebagai permukaan parametrik.
\end{eulercomment}
\begin{eulerprompt}
>function map f(a) := bisect("fr",0,2;a); ...
>t=linspace(-pi/2,pi/2,100); r=f(t);  ...
>s=linspace(pi,2pi,100)'; ...
>plot3d(r*cos(t)*sin(s),r*cos(t)*cos(s),r*sin(t), ...
>>hue,<frame,color=red,zoom=4,amb=0,max=0.7,grid=12,height=50°):
\end{eulerprompt}
\eulerimg{17}{images/EMT4Plot3D_Indonesia-051.png}
\begin{eulercomment}
Berikut adalah plot 3D dari gambar di atas yang diputar mengelilingi sumbu z. Kita mendefinisikan fungsi yang
menggambarkan objek tersebut.
\end{eulercomment}
\begin{eulerprompt}
>function f(x,y,z) ...
\end{eulerprompt}
\begin{eulerudf}
  r=x^2+y^2;
  return (r+z^2-1)^3-r*z^3;
   endfunction
\end{eulerudf}
\begin{eulerprompt}
>plot3d("f(x,y,z)", ...
>xmin=0,xmax=1.2,ymin=-1.2,ymax=1.2,zmin=-1.2,zmax=1.4, ...
>implicit=1,angle=-30°,zoom=2.5,n=[10,100,60],>anaglyph):
\end{eulerprompt}
\eulerimg{17}{images/EMT4Plot3D_Indonesia-052.png}
\eulerheading{Plot 3D Khusus}
\begin{eulercomment}
Fungsi plot3d memang bagus, tetapi tidak memenuhi semua kebutuhan. Selain rutin-rutin yang lebih dasar, dimungkinkan
mendapatkan plot berbingkai dari objek apa pun yang Anda inginkan.

Meskipun Euler bukan program 3D, Euler dapat menggabungkan beberapa objek dasar. Kita mencoba memvisualisasikan
paraboloid dan garis singgungnya.
\end{eulercomment}
\begin{eulerprompt}
>function myplot ...
\end{eulerprompt}
\begin{eulerudf}
    y=-1:0.01:1; x=(-1:0.01:1)';
    plot3d(x,y,0.2*(x-0.1)/2,<scale,<frame,>hue, ..
      hues=0.5,>contour,color=orange);
    h=holding(1);
    plot3d(x,y,(x^2+y^2)/2,<scale,<frame,>contour,>hue);
    holding(h);
  endfunction
\end{eulerudf}
\begin{eulercomment}
Sekarang framedplot() menyediakan bingkai, dan mengatur tampilan.
\end{eulercomment}
\begin{eulerprompt}
>framedplot("myplot",[-1,1,-1,1,0,1],height=0,angle=-30°, ...
>  center=[0,0,-0.7],zoom=3):
\end{eulerprompt}
\eulerimg{17}{images/EMT4Plot3D_Indonesia-053.png}
\begin{eulercomment}
Dengan cara yang sama, Anda dapat menggambar bidang kontur secara manual. Perhatikan bahwa plot3d() secara default
mengatur jendela menjadi fullwindow(), tetapi plotcontourplane() mengasumsikan hal itu.
\end{eulercomment}
\begin{eulerprompt}
>x=-1:0.02:1.1; y=x'; z=x^2-y^4;
>function myplot (x,y,z) ...
\end{eulerprompt}
\begin{eulerudf}
    zoom(2);
    wi=fullwindow();
    plotcontourplane(x,y,z,level="auto",<scale);
    plot3d(x,y,z,>hue,<scale,>add,color=white,level="thin");
    window(wi);
    reset();
  endfunction
\end{eulerudf}
\begin{eulerprompt}
>myplot(x,y,z):
\end{eulerprompt}
\eulerimg{13}{images/EMT4Plot3D_Indonesia-054.png}
\eulerheading{Animasi}
\begin{eulercomment}
Euler dapat menggunakan frame untuk menghitung animasi terlebih dahulu.

Salah satu fungsi yang memanfaatkan teknik ini adalah rotate. Fungsi ini dapat mengubah sudut pandang dan menggambar
ulang plot 3D. Fungsi ini memanggil addpage() untuk setiap plot baru. Terakhir fungsi ini menganimasikan plot-plot
tersebut.

Silakan pelajari kode sumber rotate untuk melihat lebih banyak rincian.
\end{eulercomment}
\begin{eulerprompt}
>function testplot () := plot3d("x^2+y^3"); ...
>rotate("testplot"); testplot():
\end{eulerprompt}
\eulerimg{13}{images/EMT4Plot3D_Indonesia-055.png}
\eulerheading{Menggambar Povray}
\begin{eulercomment}
Dengan bantuan file Euler povray.e, Euler dapat menghasilkan file Povray. Hasilnya sangat enak dipandang.

Anda perlu memasang Povray (32bit atau 64bit) dari http://www.povray.org/, dan menaruh subdirektori "bin" Povray ke
dalam environment path, atau mengatur variabel "defaultpovray" dengan path lengkap yang menunjuk ke "pvengine.exe".

Antarmuka Povray Euler menghasilkan file Povray di direktori home pengguna, dan memanggil Povray untuk mem-parse
file-file tersebut. Nama file defaultnya adalah current.pov, dan direktori defaultnya adalah eulerhome(), biasanya
c:\textbackslash{}Users\textbackslash{}Username\textbackslash{}Euler. Povray menghasilkan file PNG, yang dapat dimuat Euler ke dalam notebook. Untuk membersihkan
file-file ini, gunakan povclear().

Fungsi pov3d sejalan dengan plot3d. Fungsi ini dapat menghasilkan grafik fungsi f(x,y), atau permukaan dengan
koordinat X,Y,Z dalam matriks, termasuk garis level opsional. Fungsi ini menjalankan raytracer secara otomatis, dan
memuat adegan (scene) ke dalam notebook Euler.

Selain pov3d(), ada banyak fungsi yang menghasilkan objek Povray. Fungsi-fungsi ini mengembalikan string yang berisi
kode Povray untuk objek tersebut. Untuk menggunakan fungsi-fungsi ini, mulailah file Povray dengan povstart().
Kemudian gunakan writeln(...) untuk menulis objek ke file scene. Terakhir, akhiri file dengan povend(). Secara
default, raytracer akan berjalan, dan PNG akan disisipkan ke notebook Euler.

Fungsi-fungsi objek memiliki parameter bernama "look", yang memerlukan string berisi kode Povray untuk tekstur dan
finishing objek. Fungsi povlook() dapat digunakan untuk menghasilkan string ini. Fungsi ini memiliki parameter untuk
warna, transparansi, Phong Shading, dan sebagainya.

Perhatikan bahwa dunia Povray memiliki sistem koordinat lain. Antarmuka ini menerjemahkan semua koordinat ke sistem
Povray. Jadi Anda tetap dapat berpikir dalam sistem koordinat Euler dengan z mengarah vertikal ke atas, dan sumbu
x,y,z menurut kaidah tangan kanan.\\
Anda perlu memuat file povray.
\end{eulercomment}
\begin{eulerprompt}
>load povray;
\end{eulerprompt}
\begin{eulercomment}
Pastikan direktori bin Povray ada dalam path. Jika tidak, ubahlah variabel berikut sehingga memuat path ke file
eksekusi povray.
\end{eulercomment}
\begin{eulerprompt}
>defaultpovray="C:\(\backslash\)Program Files\(\backslash\)POV-Ray\(\backslash\)v3.7\(\backslash\)bin\(\backslash\)pvengine.exe"
\end{eulerprompt}
\begin{euleroutput}
  C:\(\backslash\)Program Files\(\backslash\)POV-Ray\(\backslash\)v3.7\(\backslash\)bin\(\backslash\)pvengine.exe
\end{euleroutput}
\begin{eulercomment}
Untuk kesan pertama, kita menggambar fungsi sederhana. Perintah berikut menghasilkan file povray di direktori
pengguna Anda, dan menjalankan Povray untuk me-ray-trace file ini.

Jika Anda menjalankan perintah berikut, GUI Povray akan terbuka, menjalankan file, dan menutup secara otomatis.
Karena alasan keamanan, Anda akan ditanya apakah ingin mengizinkan file exe berjalan. Anda dapat menekan cancel
untuk menghentikan pertanyaan berikutnya. Anda mungkin perlu menekan OK pada jendela Povray untuk mengonfirmasi
dialog awal Povray.
\end{eulercomment}
\begin{eulerprompt}
>plot3d("x^2+y^2",zoom=2):
\end{eulerprompt}
\eulerimg{13}{images/EMT4Plot3D_Indonesia-056.png}
\begin{eulerprompt}
>pov3d("x^2+y^2",zoom=3);
\end{eulerprompt}
\begin{euleroutput}
  exec:
      return _exec(program,param,dir,print,hidden,wait);
  povray:
      exec(program,params,defaulthome);
  Try "trace errors" to inspect local variables after errors.
  pov3d:
      if povray then povray(currentfile,w,h,w/h); endif;
\end{euleroutput}
\begin{eulercomment}
Kita dapat membuat fungsi transparan dan menambahkan finishing lain. Kita juga dapat menambahkan garis level pada
plot fungsi.
\end{eulercomment}
\begin{eulerprompt}
>pov3d("x^2+y^3",axiscolor=red,angle=-45°,>anaglyph, ...
>  look=povlook(cyan,0.2),level=-1:0.5:1,zoom=3.8);
\end{eulerprompt}
\begin{euleroutput}
  exec:
      return _exec(program,param,dir,print,hidden,wait);
  povray:
      exec(program,params,defaulthome);
  Try "trace errors" to inspect local variables after errors.
  pov3d:
      if povray then povray(currentfile,w,h,w/h); endif;
\end{euleroutput}
\begin{eulercomment}
Terkadang perlu mencegah penskalaan fungsi, dan menskalakan fungsi secara manual.

Kita menggambar himpunan titik pada bidang kompleks, di mana hasil kali jarak ke 1 dan -1 sama dengan 1.
\end{eulercomment}
\begin{eulerprompt}
>pov3d("((x-1)^2+y^2)*((x+1)^2+y^2)/40",r=2, ...
>  angle=-120°,level=1/40,dlevel=0.005,light=[-1,1,1],height=10°,n=50, ...
>  <fscale,zoom=3.8);
\end{eulerprompt}
\begin{euleroutput}
  exec:
      return _exec(program,param,dir,print,hidden,wait);
  povray:
      exec(program,params,defaulthome);
  Try "trace errors" to inspect local variables after errors.
  pov3d:
      if povray then povray(currentfile,w,h,w/h); endif;
\end{euleroutput}
\eulerheading{Menggambar dengan Koordinat}
\begin{eulercomment}
Sebagai ganti fungsi, kita dapat menggambar dengan koordinat. Seperti pada plot3d, kita memerlukan tiga matriks
untuk mendefinisikan objek.

Pada contoh ini kita memutar suatu fungsi mengelilingi sumbu z.
\end{eulercomment}
\begin{eulerprompt}
>function f(x) := x^3-x+1; ...
>x=-1:0.01:1; t=linspace(0,2pi,50)'; ...
>Z=x; X=cos(t)*f(x); Y=sin(t)*f(x); ...
>pov3d(X,Y,Z,angle=40°,look=povlook(red,0.1),height=50°,axis=0,zoom=4,light=[10,5,15]);
\end{eulerprompt}
\begin{euleroutput}
  Command was not allowed!
  exec:
      return _exec(program,param,dir,print,hidden,wait);
  povray:
      exec(program,params,defaulthome);
  Try "trace errors" to inspect local variables after errors.
  pov3d:
      if povray then povray(currentfile,w,h,w/h); endif;
\end{euleroutput}
\begin{eulercomment}
Pada contoh berikut, kita menggambar gelombang teredam. Kita membangkitkan gelombang dengan bahasa matriks Euler.

Kita juga menunjukkan bagaimana objek tambahan dapat ditambahkan ke scene pov3d. Untuk pembuatan objek, lihat
contoh-contoh berikut. Perhatikan bahwa plot3d menskalakan plot, sehingga muat dalam kubus satuan.
\end{eulercomment}
\begin{eulerprompt}
>r=linspace(0,1,80); phi=linspace(0,2pi,80)'; ...
>x=r*cos(phi); y=r*sin(phi); z=exp(-5*r)*cos(8*pi*r)/3;  ...
>pov3d(x,y,z,zoom=6,axis=0,height=30°,add=povsphere([0.5,0,0.25],0.15,povlook(red)), ...
>  w=500,h=300);
\end{eulerprompt}
\begin{euleroutput}
  Command was not allowed!
  exec:
      return _exec(program,param,dir,print,hidden,wait);
  povray:
      exec(program,params,defaulthome);
  Try "trace errors" to inspect local variables after errors.
  pov3d:
      if povray then povray(currentfile,w,h,w/h); endif;
\end{euleroutput}
\begin{eulercomment}
Dengan metode shading lanjutan Povray, sangat sedikit titik dapat menghasilkan permukaan yang sangat mulus. Hanya
pada batas dan bayangan trik ini mungkin terlihat.

Untuk ini, kita perlu menambahkan vektor normal pada setiap titik matriks.
\end{eulercomment}
\begin{eulerprompt}
>Z &= x^2*y^3
\end{eulerprompt}
\begin{euleroutput}
  
                                   2  3
                                  x  y
  
\end{euleroutput}
\begin{eulercomment}
Persamaan permukaannya adalah [x,y,Z]. Kita menghitung dua turunan terhadap x dan y dari persamaan ini dan mengambil
hasil kali silangnya sebagai normal.
\end{eulercomment}
\begin{eulerprompt}
>dx &= diff([x,y,Z],x); dy &= diff([x,y,Z],y);
\end{eulerprompt}
\begin{eulercomment}
Kita definisikan normal sebagai hasil kali silang turunan-turunan ini, dan mendefinisikan fungsi-fungsi koordinat.
\end{eulercomment}
\begin{eulerprompt}
>N &= crossproduct(dx,dy); NX &= N[1]; NY &= N[2]; NZ &= N[3]; N,
\end{eulerprompt}
\begin{euleroutput}
  
                                 3       2  2
                         [- 2 x y , - 3 x  y , 1]
  
\end{euleroutput}
\begin{eulercomment}
Kita hanya menggunakan 25 titik.
\end{eulercomment}
\begin{eulerprompt}
>x=-1:0.5:1; y=x';
>pov3d(x,y,Z(x,y),angle=10°, ...
>  xv=NX(x,y),yv=NY(x,y),zv=NZ(x,y),<shadow);
\end{eulerprompt}
\begin{euleroutput}
  exec:
      return _exec(program,param,dir,print,hidden,wait);
  povray:
      exec(program,params,defaulthome);
  Try "trace errors" to inspect local variables after errors.
  pov3d:
      if povray then povray(currentfile,w,h,w/h); endif;
\end{euleroutput}
\begin{eulercomment}
Berikut adalah simpul Trefoil yang dibuat oleh A. Busser di Povray. Ada versi yang lebih baik dari ini pada
contoh-contoh.

Lihat: Examples\textbackslash{}Trefoil Knot \textbar{} Trefoil Knot

Agar tampilannya bagus dengan titik yang tidak terlalu banyak, kita tambahkan vektor normal di sini. Kita
menggunakan Maxima untuk menghitung normal. Pertama, tiga fungsi untuk koordinat sebagai ekspresi simbolik.
\end{eulercomment}
\begin{eulerprompt}
>X &= ((4+sin(3*y))+cos(x))*cos(2*y); ...
>Y &= ((4+sin(3*y))+cos(x))*sin(2*y); ...
>Z &= sin(x)+2*cos(3*y);
\end{eulerprompt}
\begin{eulercomment}
Kemudian dua vektor turunan terhadap x dan y.
\end{eulercomment}
\begin{eulerprompt}
>dx &= diff([X,Y,Z],x); dy &= diff([X,Y,Z],y);
\end{eulerprompt}
\begin{eulercomment}
Sekarang normal, yang merupakan hasil kali silang dari kedua turunan.
\end{eulercomment}
\begin{eulerprompt}
>dn &= crossproduct(dx,dy);
\end{eulerprompt}
\begin{eulercomment}
Sekarang kita evaluasi semua ini secara numerik.
\end{eulercomment}
\begin{eulerprompt}
>x:=linspace(-%pi,%pi,40); y:=linspace(-%pi,%pi,100)';
\end{eulerprompt}
\begin{eulercomment}
Vektor normal adalah evaluasi dari ekspresi simbolik dn[i] untuk i=1,2,3. Sintaksnya adalah \&"ekspresi"(parameter).
Ini adalah alternatif dari metode pada contoh sebelumnya, di mana kita mendefinisikan terlebih dahulu ekspresi
simbolik NX, NY, NZ.
\end{eulercomment}
\begin{eulerprompt}
>pov3d(X(x,y),Y(x,y),Z(x,y),>anaglyph,axis=0,zoom=5,w=450,h=350, ...
>  <shadow,look=povlook(blue), ...
>  xv=&"dn[1]"(x,y), yv=&"dn[2]"(x,y), zv=&"dn[3]"(x,y));
\end{eulerprompt}
\begin{euleroutput}
  Command was not allowed!
  exec:
      return _exec(program,param,dir,print,hidden,wait);
  povray:
      exec(program,params,defaulthome);
  Try "trace errors" to inspect local variables after errors.
  pov3d:
      if povray then povray(currentfile,w,h,w/h); endif;
\end{euleroutput}
\begin{eulercomment}
Kita juga dapat membangkitkan grid dalam 3D.
\end{eulercomment}
\begin{eulerprompt}
>povstart(zoom=4); ...
>x=-1:0.5:1; r=1-(x+1)^2/6; ...
>t=(0°:30°:360°)'; y=r*cos(t); z=r*sin(t); ...
>writeln(povgrid(x,y,z,d=0.02,dballs=0.05)); ...
>povend();
\end{eulerprompt}
\begin{euleroutput}
  exec:
      return _exec(program,param,dir,print,hidden,wait);
  povray:
      exec(program,params,defaulthome);
  Try "trace errors" to inspect local variables after errors.
  povend:
      povray(file,w,h,aspect,exit); 
\end{euleroutput}
\begin{eulercomment}
Dengan povgrid(), kurva dimungkinkan.
\end{eulercomment}
\begin{eulerprompt}
>povstart(center=[0,0,1],zoom=3.6); ...
>t=linspace(0,2,1000); r=exp(-t); ...
>x=cos(2*pi*10*t)*r; y=sin(2*pi*10*t)*r; z=t; ...
>writeln(povgrid(x,y,z,povlook(red))); ...
>writeAxis(0,2,axis=3); ...
>povend();
\end{eulerprompt}
\begin{euleroutput}
  Command was not allowed!
  exec:
      return _exec(program,param,dir,print,hidden,wait);
  povray:
      exec(program,params,defaulthome);
  Try "trace errors" to inspect local variables after errors.
  povend:
      povray(file,w,h,aspect,exit); 
\end{euleroutput}
\eulerheading{Objek Povray}
\begin{eulercomment}
Di atas, kita menggunakan pov3d untuk menggambar permukaan. Antarmuka povray di Euler juga dapat menghasilkan objek
Povray. Objek-objek ini disimpan sebagai string di Euler, dan perlu ditulis ke file Povray.

Kita mulai keluaran dengan povstart().
\end{eulercomment}
\begin{eulerprompt}
>povstart(zoom=4);
\end{eulerprompt}
\begin{eulercomment}
Pertama kita definisikan tiga silinder, dan menyimpannya dalam string di Euler.

Fungsi povx() dan sebagainya hanya mengembalikan vektor [1,0,0], yang dapat digunakan sebagai gantinya.
\end{eulercomment}
\begin{eulerprompt}
>c1=povcylinder(-povx,povx,1,povlook(red)); ...
>c2=povcylinder(-povy,povy,1,povlook(yellow)); ...
>c3=povcylinder(-povz,povz,1,povlook(blue)); ...
\end{eulerprompt}
\begin{eulercomment}
String tersebut berisi kode Povray, yang tidak perlu kita pahami pada tahap ini.
\end{eulercomment}
\begin{eulerprompt}
>c2
\end{eulerprompt}
\begin{euleroutput}
  cylinder \{ <0,0,-1>, <0,0,1>, 1
   texture \{ pigment \{ color rgb <0.941176,0.941176,0.392157> \}  \} 
   finish \{ ambient 0.2 \} 
   \}
\end{euleroutput}
\begin{eulercomment}
Seperti yang Anda lihat, kita menambahkan tekstur pada objek dengan tiga warna berbeda.

Hal itu dilakukan oleh povlook(), yang mengembalikan string berisi kode Povray yang relevan. Kita dapat menggunakan
warna default Euler, atau mendefinisikan warna sendiri. Kita juga dapat menambahkan transparansi, atau mengubah
cahaya sekitar (ambient).
\end{eulercomment}
\begin{eulerprompt}
>povlook(rgb(0.1,0.2,0.3),0.1,0.5)
\end{eulerprompt}
\begin{euleroutput}
   texture \{ pigment \{ color rgbf <0.101961,0.2,0.301961,0.1> \}  \} 
   finish \{ ambient 0.5 \} 
  
\end{euleroutput}
\begin{eulercomment}
Sekarang kita definisikan objek irisan (intersection), dan menulis hasilnya ke file.
\end{eulercomment}
\begin{eulerprompt}
>writeln(povintersection([c1,c2,c3]));
\end{eulerprompt}
\begin{eulercomment}
Irisan tiga silinder sulit divisualisasikan, jika Anda belum pernah melihatnya sebelumnya.
\end{eulercomment}
\begin{eulerprompt}
>povend;
\end{eulerprompt}
\begin{euleroutput}
  Command was not allowed!
  exec:
      return _exec(program,param,dir,print,hidden,wait);
  povray:
      exec(program,params,defaulthome);
  Try "trace errors" to inspect local variables after errors.
  povend:
      povray(file,w,h,aspect,exit); 
\end{euleroutput}
\begin{eulercomment}
Fungsi-fungsi berikut membangkitkan fraktal secara rekursif.

Fungsi pertama menunjukkan bagaimana Euler menangani objek Povray sederhana. Fungsi povbox() mengembalikan string
yang berisi koordinat kotak, tekstur, dan finishing.
\end{eulercomment}
\begin{eulerprompt}
>function onebox(x,y,z,d) := povbox([x,y,z],[x+d,y+d,z+d],povlook());
>function fractal (x,y,z,h,n) ...
\end{eulerprompt}
\begin{eulerudf}
   if n==1 then writeln(onebox(x,y,z,h));
   else
     h=h/3;
     fractal(x,y,z,h,n-1);
     fractal(x+2*h,y,z,h,n-1);
     fractal(x,y+2*h,z,h,n-1);
     fractal(x,y,z+2*h,h,n-1);
     fractal(x+2*h,y+2*h,z,h,n-1);
     fractal(x+2*h,y,z+2*h,h,n-1);
     fractal(x,y+2*h,z+2*h,h,n-1);
     fractal(x+2*h,y+2*h,z+2*h,h,n-1);
     fractal(x+h,y+h,z+h,h,n-1);
   endif;
  endfunction
\end{eulerudf}
\begin{eulerprompt}
>povstart(fade=10,<shadow);
>fractal(-1,-1,-1,2,4);
>povend();
\end{eulerprompt}
\begin{euleroutput}
  Command was not allowed!
  exec:
      return _exec(program,param,dir,print,hidden,wait);
  povray:
      exec(program,params,defaulthome);
  Try "trace errors" to inspect local variables after errors.
  povend:
      povray(file,w,h,aspect,exit); 
\end{euleroutput}
\begin{eulercomment}
Selisih (difference) memungkinkan memotong satu objek dari objek lain. Seperti irisan, ini merupakan bagian dari
objek CSG Povray.
\end{eulercomment}
\begin{eulerprompt}
>povstart(light=[5,-5,5],fade=10);
\end{eulerprompt}
\begin{eulercomment}
Untuk demonstrasi ini, kita mendefinisikan objek di Povray, sebagai ganti menggunakan string di Euler. Definisi
langsung ditulis ke file.

Koordinat kotak -1 berarti [-1,-1,-1].
\end{eulercomment}
\begin{eulerprompt}
>povdefine("mycube",povbox(-1,1));
\end{eulerprompt}
\begin{eulercomment}
Kita dapat menggunakan objek ini pada povobject(), yang mengembalikan string seperti biasa.
\end{eulercomment}
\begin{eulerprompt}
>c1=povobject("mycube",povlook(red));
\end{eulerprompt}
\begin{eulercomment}
Kita membuat kubus kedua, lalu memutar dan menskalakannya sedikit.
\end{eulercomment}
\begin{eulerprompt}
>c2=povobject("mycube",povlook(yellow),translate=[1,1,1], ...
>  rotate=xrotate(10°)+yrotate(10°), scale=1.2);
\end{eulerprompt}
\begin{eulercomment}
Kemudian kita ambil selisih kedua objek.
\end{eulercomment}
\begin{eulerprompt}
>writeln(povdifference(c1,c2));
\end{eulerprompt}
\begin{eulercomment}
Sekarang tambahkan tiga sumbu.
\end{eulercomment}
\begin{eulerprompt}
>writeAxis(-1.2,1.2,axis=1); ...
>writeAxis(-1.2,1.2,axis=2); ...
>writeAxis(-1.2,1.2,axis=4); ...
>povend();
\end{eulerprompt}
\begin{euleroutput}
  Command was not allowed!
  exec:
      return _exec(program,param,dir,print,hidden,wait);
  povray:
      exec(program,params,defaulthome);
  Try "trace errors" to inspect local variables after errors.
  povend:
      povray(file,w,h,aspect,exit); 
\end{euleroutput}
\eulerheading{Fungsi Implisit}
\begin{eulercomment}
Povray dapat menggambar himpunan di mana f(x,y,z)=0, seperti parameter implicit pada plot3d. Namun hasilnya tampak
jauh lebih bagus.

Sintaks untuk fungsinya sedikit berbeda. Anda tidak dapat menggunakan keluaran ekspresi Maxima atau Euler.

\end{eulercomment}
\begin{eulerformula}
\[
((x^2+y^2-c^2)^2+(z^2-1)^2)*((y^2+z^2-c^2)^2+(x^2-1)^2)*((z^2+x^2-c^2)^2+(y^2-1)^2)=d
\]
\end{eulerformula}
\begin{eulerprompt}
>povstart(angle=70°,height=50°,zoom=4);
>c=0.1; d=0.1; ...
>writeln(povsurface("(pow(pow(x,2)+pow(y,2)-pow(c,2),2)+pow(pow(z,2)-1,2))*(pow(pow(y,2)+pow(z,2)-pow(c,2),2)+pow(pow(x,2)-1,2))*(pow(pow(z,2)+pow(x,2)-pow(c,2),2)+pow(pow(y,2)-1,2))-d",povlook(red))); ...
>povend();
\end{eulerprompt}
\begin{euleroutput}
  exec:
      return _exec(program,param,dir,print,hidden,wait);
  povray:
      exec(program,params,defaulthome);
  Try "trace errors" to inspect local variables after errors.
  povend:
      povray(file,w,h,aspect,exit); 
\end{euleroutput}
\begin{eulerprompt}
>povstart(angle=25°,height=10°); 
>writeln(povsurface("pow(x,2)+pow(y,2)*pow(z,2)-1",povlook(blue),povbox(-2,2,"")));
>povend();
\end{eulerprompt}
\begin{euleroutput}
  Command was not allowed!
  exec:
      return _exec(program,param,dir,print,hidden,wait);
  povray:
      exec(program,params,defaulthome);
  Try "trace errors" to inspect local variables after errors.
  povend:
      povray(file,w,h,aspect,exit); 
\end{euleroutput}
\begin{eulerprompt}
>povstart(angle=70°,height=50°,zoom=4);
\end{eulerprompt}
\begin{eulercomment}
Buat permukaan implisit. Perhatikan sintaks yang berbeda pada ekspresi.
\end{eulercomment}
\begin{eulerprompt}
>writeln(povsurface("pow(x,2)*y-pow(y,3)-pow(z,2)",povlook(green))); ...
>writeAxes(); ...
>povend();
\end{eulerprompt}
\begin{euleroutput}
  Command was not allowed!
  exec:
      return _exec(program,param,dir,print,hidden,wait);
  povray:
      exec(program,params,defaulthome);
  Try "trace errors" to inspect local variables after errors.
  povend:
      povray(file,w,h,aspect,exit); 
\end{euleroutput}
\eulerheading{Objek Mesh}
\begin{eulercomment}
Pada contoh ini, kita menunjukkan cara membuat objek mesh, dan menggambarnya dengan informasi tambahan.

Kita ingin memaksimalkan xy dengan syarat x+y=1 dan mendemonstrasikan persinggungan garis-garis level.
\end{eulercomment}
\begin{eulerprompt}
>povstart(angle=-10°,center=[0.5,0.5,0.5],zoom=7);
\end{eulerprompt}
\begin{eulercomment}
Kita tidak dapat menyimpan objek dalam string seperti sebelumnya, karena terlalu besar. Jadi kita mendefinisikan
objek dalam file Povray menggunakan #declare. Fungsi povtriangle() melakukan ini secara otomatis. Fungsi ini dapat
menerima vektor normal seperti pov3d().

Berikut mendefinisikan objek mesh, dan langsung menuliskannya ke file.
\end{eulercomment}
\begin{eulerprompt}
>x=0:0.02:1; y=x'; z=x*y; vx=-y; vy=-x; vz=1;
>mesh=povtriangles(x,y,z,"",vx,vy,vz);
\end{eulerprompt}
\begin{eulercomment}
Sekarang kita definisikan dua cakram, yang akan diiriskan dengan permukaan.
\end{eulercomment}
\begin{eulerprompt}
>cl=povdisc([0.5,0.5,0],[1,1,0],2); ...
>ll=povdisc([0,0,1/4],[0,0,1],2);
\end{eulerprompt}
\begin{eulercomment}
Tulis permukaan dikurangi kedua cakram.
\end{eulercomment}
\begin{eulerprompt}
>writeln(povdifference(mesh,povunion([cl,ll]),povlook(green)));
\end{eulerprompt}
\begin{eulercomment}
Tulis kedua irisan.
\end{eulercomment}
\begin{eulerprompt}
>writeln(povintersection([mesh,cl],povlook(red))); ...
>writeln(povintersection([mesh,ll],povlook(gray)));
\end{eulerprompt}
\begin{eulercomment}
Tulis sebuah titik pada maksimum.
\end{eulercomment}
\begin{eulerprompt}
>writeln(povpoint([1/2,1/2,1/4],povlook(gray),size=2*defaultpointsize));
\end{eulerprompt}
\begin{eulercomment}
Tambahkan sumbu dan selesaikan.
\end{eulercomment}
\begin{eulerprompt}
>writeAxes(0,1,0,1,0,1,d=0.015); ...
>povend();
\end{eulerprompt}
\begin{euleroutput}
  Command was not allowed!
  exec:
      return _exec(program,param,dir,print,hidden,wait);
  povray:
      exec(program,params,defaulthome);
  Try "trace errors" to inspect local variables after errors.
  povend:
      povray(file,w,h,aspect,exit); 
\end{euleroutput}
\eulerheading{Anaglyph di Povray}
\begin{eulercomment}
Untuk menghasilkan anaglyph untuk kacamata merah/cyan, Povray harus dijalankan dua kali dari posisi kamera yang
berbeda. Ini menghasilkan dua file Povray dan dua file PNG, yang dimuat dengan fungsi loadanaglyph().

Tentu saja, Anda memerlukan kacamata merah/cyan untuk melihat contoh-contoh berikut dengan benar.

Fungsi pov3d() memiliki saklar sederhana untuk menghasilkan anaglyph.
\end{eulercomment}
\begin{eulerprompt}
>pov3d("-exp(-x^2-y^2)/2",r=2,height=45°,>anaglyph, ...
>  center=[0,0,0.5],zoom=3.5);
\end{eulerprompt}
\begin{euleroutput}
  exec:
      return _exec(program,param,dir,print,hidden,wait);
  povray:
      exec(program,params,defaulthome);
  Try "trace errors" to inspect local variables after errors.
  pov3d:
      if povray then povray(currentfile,w,h,w/h); endif;
\end{euleroutput}
\begin{eulercomment}
Jika Anda membuat scene dengan objek, Anda perlu menaruh pembuatan scene ke dalam sebuah fungsi, dan menjalankannya
dua kali dengan nilai berbeda untuk parameter anaglyph.
\end{eulercomment}
\begin{eulerprompt}
>function myscene ...
\end{eulerprompt}
\begin{eulerudf}
    s=povsphere(povc,1);
    cl=povcylinder(-povz,povz,0.5);
    clx=povobject(cl,rotate=xrotate(90°));
    cly=povobject(cl,rotate=yrotate(90°));
    c=povbox([-1,-1,0],1);
    un=povunion([cl,clx,cly,c]);
    obj=povdifference(s,un,povlook(red));
    writeln(obj);
    writeAxes();
  endfunction
\end{eulerudf}
\begin{eulercomment}
Fungsi povanaglyph() melakukan semua ini. Parameternya seperti gabungan povstart() dan povend().
\end{eulercomment}
\begin{eulerprompt}
>povanaglyph("myscene",zoom=4.5);
\end{eulerprompt}
\begin{euleroutput}
  Command was not allowed!
  exec:
      return _exec(program,param,dir,print,hidden,wait);
  povray:
      exec(program,params,defaulthome);
  Try "trace errors" to inspect local variables after errors.
  povanaglyph:
      povray(currentfile,w,h,aspect,exit); 
\end{euleroutput}
\eulerheading{Mendefinisikan Objek Sendiri}
\begin{eulercomment}
Antarmuka povray Euler berisi banyak objek. Tetapi Anda tidak dibatasi pada objek-objek tersebut. Anda dapat membuat
objek sendiri, yang menggabungkan objek lain, atau berupa objek yang sepenuhnya baru.

Kita mendemonstrasikan torus. Perintah Povray untuknya adalah "torus". Jadi kita mengembalikan string dengan
perintah ini beserta parameternya. Perhatikan bahwa torus selalu berpusat di titik asal.
\end{eulercomment}
\begin{eulerprompt}
>function povdonat (r1,r2,look="") ...
\end{eulerprompt}
\begin{eulerudf}
    return "torus \{"+r1+","+r2+look+"\}";
  endfunction
\end{eulerudf}
\begin{eulercomment}
Berikut torus pertama kita.
\end{eulercomment}
\begin{eulerprompt}
>t1=povdonat(0.8,0.2)
\end{eulerprompt}
\begin{euleroutput}
  torus \{0.8,0.2\}
\end{euleroutput}
\begin{eulercomment}
Mari kita gunakan objek ini untuk membuat torus kedua, yang digeser dan diputar.
\end{eulercomment}
\begin{eulerprompt}
>t2=povobject(t1,rotate=xrotate(90°),translate=[0.8,0,0])
\end{eulerprompt}
\begin{euleroutput}
  object \{ torus \{0.8,0.2\}
   rotate 90 *x 
   translate <0.8,0,0>
   \}
\end{euleroutput}
\begin{eulercomment}
Sekarang kita tempatkan objek-objek ini ke dalam scene. Untuk tampilannya, kita menggunakan Phong Shading.
\end{eulercomment}
\begin{eulerprompt}
>povstart(center=[0.4,0,0],angle=0°,zoom=3.8,aspect=1.5); ...
>writeln(povobject(t1,povlook(green,phong=1))); ...
>writeln(povobject(t2,povlook(green,phong=1))); ...
\end{eulerprompt}
\begin{eulerttcomment}
 >povend();
\end{eulerttcomment}
\begin{eulercomment}
memanggil program Povray. Namun, jika terjadi galat, galatnya tidak ditampilkan. Karena itu Anda sebaiknya
menggunakan

\end{eulercomment}
\begin{eulerttcomment}
 >povend(<exit);
\end{eulerttcomment}
\begin{eulercomment}

jika ada yang tidak berfungsi. Ini akan membiarkan jendela Povray tetap terbuka.
\end{eulercomment}
\begin{eulerprompt}
>povend(h=320,w=480);
\end{eulerprompt}
\begin{euleroutput}
  Command was not allowed!
  exec:
      return _exec(program,param,dir,print,hidden,wait);
  povray:
      exec(program,params,defaulthome);
  Try "trace errors" to inspect local variables after errors.
  povend:
      povray(file,w,h,aspect,exit); 
\end{euleroutput}
\begin{eulercomment}
Berikut contoh yang lebih rumit. Kita menyelesaikan

\end{eulercomment}
\begin{eulerformula}
\[
Ax \le b, \quad x \ge 0, \quad c.x \to \text{Max.}
\]
\end{eulerformula}
\begin{eulercomment}
dan menampilkan titik-titik layak (feasible) dan titik optimum dalam plot 3D.
\end{eulercomment}
\begin{eulerprompt}
>A=[10,8,4;5,6,8;6,3,2;9,5,6];
>b=[10,10,10,10]';
>c=[1,1,1];
\end{eulerprompt}
\begin{eulercomment}
Pertama, mari kita periksa apakah contoh ini memiliki solusi.
\end{eulercomment}
\begin{eulerprompt}
>x=simplex(A,b,c,>max,>check)'
\end{eulerprompt}
\begin{euleroutput}
  [0,  1,  0.5]
\end{euleroutput}
\begin{eulercomment}
Ya, ada.

Selanjutnya kita definisikan dua objek. Yang pertama adalah bidang

\end{eulercomment}
\begin{eulerformula}
\[
a \cdot x \le b
\]
\end{eulerformula}
\begin{eulerprompt}
>function oneplane (a,b,look="") ...
\end{eulerprompt}
\begin{eulerudf}
    return povplane(a,b,look)
  endfunction
\end{eulerudf}
\begin{eulercomment}
Kemudian kita definisikan irisan semua ruang setengah dan sebuah kubus.
\end{eulercomment}
\begin{eulerprompt}
>function adm (A, b, r, look="") ...
\end{eulerprompt}
\begin{eulerudf}
    ol=[];
    loop 1 to rows(A); ol=ol|oneplane(A[#],b[#]); end;
    ol=ol|povbox([0,0,0],[r,r,r]);
    return povintersection(ol,look);
  endfunction
\end{eulerudf}
\begin{eulercomment}
Sekarang kita dapat menggambar scene-nya.
\end{eulercomment}
\begin{eulerprompt}
>povstart(angle=120°,center=[0.5,0.5,0.5],zoom=3.5); ...
>writeln(adm(A,b,2,povlook(green,0.4))); ...
>writeAxes(0,1.3,0,1.6,0,1.5); ...
\end{eulerprompt}
\begin{eulercomment}
Berikut adalah lingkaran di sekitar titik optimum.
\end{eulercomment}
\begin{eulerprompt}
>writeln(povintersection([povsphere(x,0.5),povplane(c,c.x')], ...
>  povlook(red,0.9)));
\end{eulerprompt}
\begin{eulercomment}
Dan sebuah panah (pada teks asli tertulis "error") ke arah titik optimum.
\end{eulercomment}
\begin{eulerprompt}
>writeln(povarrow(x,c*0.5,povlook(red)));
\end{eulerprompt}
\begin{eulercomment}
Kita menambahkan teks ke layar. Teks hanyalah objek 3D. Kita perlu menempatkan dan memutarnya sesuai pandangan kita.
\end{eulercomment}
\begin{eulerprompt}
>writeln(povtext("Linear Problem",[0,0.2,1.3],size=0.05,rotate=5°)); ...
>povend();
\end{eulerprompt}
\begin{euleroutput}
  Command was not allowed!
  exec:
      return _exec(program,param,dir,print,hidden,wait);
  povray:
      exec(program,params,defaulthome);
  Try "trace errors" to inspect local variables after errors.
  povend:
      povray(file,w,h,aspect,exit); 
\end{euleroutput}
\eulerheading{Contoh Lainnya}
\begin{eulercomment}
Anda dapat menemukan beberapa contoh lain untuk Povray di Euler pada
file-file berikut.

Lihat: Examples/Dandelin Spheres\\
Lihat: Examples/Donat Math\\
Lihat: Examples/Trefoil Knot\\
Lihat: Examples/Optimization by Affine Scaling

\begin{eulercomment}
\eulerheading{contoh penyelesaian masalah menggambar kurva/plot 3D}
\begin{eulercomment}
Soal 1: Kalkulus – Visualisasi Permukaan Fungsi Dua Variabel
(Paraboloid)Soal:Gambarkan grafik fungsi dua variabel\\
\end{eulercomment}
\begin{eulerformula}
\[
f(x,y) = x^2 +  y^2
\]
\end{eulerformula}
\begin{eulercomment}
pada domain\\
\end{eulercomment}
\begin{eulerformula}
\[
x, y \in [-2, 2].
\]
\end{eulerformula}
\begin{eulercomment}
Tampilkan kurva menggunakan skema warna spektral (spectral) agar
perbedaan ketinggian z terlihat jelas.Tambahkan garis kontur (contour)
otomatis dan atur banyaknya kisi (n=100) agar kurva tampil mulus.
\end{eulercomment}
\begin{eulerprompt}
>reset;
>plot3d("x^2+y^2", r=2, n=100,  ...
>>spectral, >contour, title="Paraboloid z = x^2 + y^2"):
\end{eulerprompt}
\eulerimg{13}{images/EMT4Plot3D_Indonesia-059.png}
\begin{eulercomment}
Penjelasan Kode:plot3d("x\textasciicircum{}2+y\textasciicircum{}2", ...) memplot fungsi dua variabel z =
f(x,y).r=2 menentukan batas simetris rentang sumbu x dan y dari -2
sampai 2.n=100 menentukan densitas subinterval pembentuk jaring-jaring
permukaan (makin besar nilainya, permukaan makin mulus).\textgreater{}spectral
menambahkan gradien warna berdasarkan Ketinggian nilai z (dari biru
hingga merah).\textgreater{}contour menggambar garis-garis kontur ketinggian di
atas permukaan objek.

Soal 2: Kalkulus/Geometri – Permukaan Parametrik 3D
(Torus/Donat)Soal:Sebuah torus (objek berbentuk donat) didefinisikan
secara parametrik oleh:\\
\end{eulercomment}
\begin{eulerformula}
\[
x(u,v) = (3 + \cos u) \cos v ; y(u,v) = (3 + \cos u) \sin v
\]
\end{eulerformula}
\begin{eulerformula}
\[
z(u,v) = \sin u
\]
\end{eulerformula}
\begin{eulercomment}
dengan\\
\end{eulercomment}
\begin{eulerformula}
\[
u \in [0, 2\pi] ; v \in [0, 2\pi]
\]
\end{eulerformula}
\begin{eulercomment}
Gambarkan bentuk objek 3D tersebut menggunakan jaring-jaring kawat
(wireframe) transparan bergradien warna.
\end{eulercomment}
\begin{eulerprompt}
>reset;
>u = linspace(0, 2*pi, 60);
>v = linspace(0, 2*pi, 60)';
>X = (3 + cos(u)) * cos(v);
>Y = (3 + cos(u)) * sin(v);
>Z = sin(u);
>plot3d(X, Y, Z, >hue, >wire, scale=1.5, title="Torus Parametrik"):
\end{eulerprompt}
\eulerimg{13}{images/EMT4Plot3D_Indonesia-063.png}
\begin{eulercomment}
Penjelasan Kode:\\
\end{eulercomment}
\begin{eulerttcomment}
 (vektor baris) dan v (vektor kolom) membuat grid koordinat 2D untuk
\end{eulerttcomment}
\begin{eulercomment}
memparameterkan permukaan.\\
X, Y, dan Z dihitung menggunakan ekspresi perkalian matriks EMT
sehingga menghasilkan matriks titik koordinat 3D.\\
\textgreater{}hue mengaktifkan pencahayaan/arsiran pada jaring.\\
\textgreater{}wire menampilkan struktur kawat (mesh/wireframe) penyusun permukaan.\\
scale=1.5 memperbesar ukuran visual objek pada ruang 3D.

Soal 3: Aljabar / Geometri Ruang – Permukaan Implisit 3D (Bola \&
Ellipsoid)\\
Soal:Gambarkan grafik implisit tiga variabel untuk persamaan \\
\end{eulercomment}
\begin{eulerformula}
\[
f(x,y,z) = x^2 + y^2 + 4z^2 - 1 = 0
\]
\end{eulerformula}
\begin{eulerttcomment}
 (Ellipsoid) pada rentang 
\end{eulerttcomment}
\begin{eulerformula}
\[
r = 1.5
\]
\end{eulerformula}
\begin{eulercomment}
.Buatlah potongan permukaan sejajar bidang x-y, x-z, dan y-z
menggunakan opsi implicit=7
\end{eulercomment}
\begin{eulerprompt}
>reset;
>plot3d("x^2 + y^2 + 4*z^2 - 1", r=1.5, implicit=7, zoom=3, title="Ellipsoid Implisit"):
\end{eulerprompt}
\eulerimg{13}{images/EMT4Plot3D_Indonesia-066.png}
\begin{eulercomment}
Penjelasan Kode:implicit=7 mengeksekusi gabungan irisan bidang (1 + 2
+ 4 = 7), yaitu potongan sejajar bidang y-z (implicit=1), x-z
(implicit=2), dan x-y (implicit=4).r=1.5 menetapkan batas ruang kubus
evaluasi [-1.5, 1.5] untuk sumbu x, y, dan z.zoom=3 memperbesar jarak
pandang kamera ke objek.


Soal 4: Kalkulus Vektor – Kurva Ruang 3D (Helix / Heliks
Spiral)Soal:Suatu partikel bergerak mengikuti kurva heliks (spiral 3D)
dengan persamaan parametrik:\\
\end{eulercomment}
\begin{eulerformula}
\[
x(t) = \cos(t), \quad y(t) = \sin(t), \quad z(t) = \frac{t}{2\pi}
\]
\end{eulerformula}
\begin{eulercomment}
untuk\\
\end{eulercomment}
\begin{eulerformula}
\[
t \in [0, 6\pi].
\]
\end{eulerformula}
\begin{eulercomment}
Gambarkan lintasan kurva 3D tersebut dengan garis berwarna biru tebal.
\end{eulercomment}
\begin{eulerprompt}
>reset;
>t = linspace(0, 6*pi, 1000);
>x = cos(t);
>y = sin(t);
>z = t / (2*pi);
>plot3d(x, y, z, >wire, linewidth=3, wirecolor=blue, zoom=3, title="Heliks 3D"):
\end{eulerprompt}
\eulerimg{13}{images/EMT4Plot3D_Indonesia-069.png}
\begin{eulercomment}
Penjelasan Kode:t = linspace(0, 6*pi, 1000) membangkitkan 1000 titik
sampel variabel parameter t.plot3d(x, y, z, \textgreater{}wire, ...) menerima tiga
vektor 1D untuk menggambar kurva berupa garis tunggal di dalam ruang
3D.linewidth=3 mengatur ketebalan garis lintasan.wirecolor=blue
mengubah warna garis heliks menjadi biru.


Soal 5: Terapan / Pemodelan – Permukaan Benda PutarSoal:Sebuah kurva \\
\end{eulercomment}
\begin{eulerformula}
\[
z = f(x) = \sqrt{x}
\]
\end{eulerformula}
\begin{eulerttcomment}
 untuk 
\end{eulerttcomment}
\begin{eulerformula}
\[
x \in [0, 4]
\]
\end{eulerformula}
\begin{eulerttcomment}
 diputar mengelilingi sumbu z (menggunakan opsi rotate=2). Gambarkan
\end{eulerttcomment}
\begin{eulercomment}
permukaan benda putar yang terbentuk (paraboloid terbalik/mangkok)
dengan garis grid \\
\end{eulercomment}
\begin{eulerformula}
\[
10 \times 10.
\]
\end{eulerformula}
\begin{eulerprompt}
>reset;
>plot3d("sqrt(x)", a=0, b=4, rotate=2, grid=10, title="Benda Putar Mengelilingi Sumbu-Z"):
\end{eulerprompt}
\eulerimg{13}{images/EMT4Plot3D_Indonesia-073.png}
\begin{eulercomment}
Penjelasan Kode:plot3d("sqrt(x)", a=0, b=4, ...) mendefinisikan kurva
2D dasar\\
\end{eulercomment}
\begin{eulerformula}
\[
y = \sqrt{x}
\]
\end{eulerformula}
\begin{eulercomment}
pada interval\\
\end{eulercomment}
\begin{eulerformula}
\[
0 \le x \le 4
\]
\end{eulerformula}
\begin{eulercomment}
.rotate=2 memutar kurva tersebut 360 derajat mengelilingi sumbu z
untuk menghasilkan bangun ruang simetris putar. (Jika rotate=1 atau
rotate=true, pemutaran dilakukan mengelilingi sumbu x).grid=10 membagi
permukaan benda putar menjadi 10 segmen garis kisi horisontal dan
vertikal.
\end{eulercomment}
\end{eulernotebook}
\end{document}
